\chapter{Getting Access: Futures and Listed Options}\label{ch:m1:getting-access-futures-options}

Two traders buy the same E-mini contract at the same price at the same
instant. One pays the exchange \$1.18 for the privilege; the other, 35 cents.
Before either has any skill, the second can profit from moves a third the
size. The difference is a membership: an asset with a market price, which
can be bought, or leased by the month from someone who owns one. Futures
exchanges publish these differences openly; they are the largest single
determinant of who can run a short-horizon futures strategy at all. Options
exchanges add a further layer, the \omterm{def:m1:getting-access-futures-options:appointment}{market-maker appointment}, which buys
allocation privileges as well as lower fees. This chapter prices both, and
the contract that comes before either: the \omterm{def:m1:getting-access-equities:brokers}{clearing broker}'s.

\section{The clearing broker}

\begin{definition}[Futures commission merchant and clearing member]\label{def:m1:getting-access-futures-options:fcm}
A \emph{futures commission merchant}\index{futures commission merchant} (FCM)
is, in the United States, an entity that solicits or accepts orders for
futures, options on futures or swaps and accepts money or other assets from
customers to support them; it must be registered and a member of the
industry's self-regulatory body. A \emph{clearing member}\index{clearing
member} is a firm admitted by a clearing house to clear trades in its own
name; it guarantees its customers' trades to the clearing house and
contributes to the \omterm{def:m1:clearing-and-settlement:waterfall}{default fund} (\cref{ch:m1:clearing-and-settlement}).
\end{definition}

Every futures position is carried by a \omterm{def:m1:getting-access-futures-options:fcm}{clearing member}. A trading firm that
is not one is the customer of an FCM, which collects at least the clearing
house's margin (\cref{ch:m1:margin}) and usually more, holds the firm's
collateral in segregated accounts, charges a commission per side, and may
give the firm access to its \omterm{def:m1:getting-access-futures-options:membership}{exchange memberships}' rates through its own
programmes. A \emph{give-up}\index{give-up} works as in equities: trades
executed through one broker are given up to the clearing FCM under a
three-party agreement.

The FCM bears the firm's default risk between \omterm{def:m1:financing:margincall}{margin calls} and prices it:
by the multiple of exchange margin it demands, by intraday limits on
position and order size, by the right to liquidate. For a leveraged
strategy these terms, not the commission, are the cost of the relationship.
In stressed periods (\cref{ch:m1:stress-case-studies}) FCMs raise their
multiples and shed customers whose tail risk they dislike; a firm's access to
futures is only as secure as its FCM's appetite.

\begin{omfigure}
\begin{tikzpicture}[font=\small,
  box/.style={draw=omDef,rounded corners=2pt,minimum width=2.7cm,minimum height=0.95cm,align=center,fill=omDef!4},
  ext/.style={draw=omExo,rounded corners=2pt,minimum width=2.7cm,minimum height=0.95cm,align=center,fill=omExo!6},
  flow/.style={-{Stealth[length=2mm]},thick}]
\node[box] (f) at (0,0) {Trading firm};
\node[box] (fcm) at (5.4,0) {Clearing FCM\\(margin, limits)};
\node[ext] (ch) at (10.8,0) {Clearing house};
\node[ext] (ex) at (10.8,2.0) {Exchange};
\node[box] (eb) at (5.4,2.0) {Executing broker\\or own membership};
\draw[flow,omDef] (f.north) |- node[pos=0.75,above,font=\footnotesize]{orders} (eb.west);
\draw[flow,omDef] (eb) -- (ex);
\draw[flow,omThm] (ex) -- node[right,font=\footnotesize]{trades} (ch);
\draw[flow,omThm] ([yshift=4pt]ch.west) -- node[above,font=\footnotesize]{margin calls} ([yshift=4pt]fcm.east);
\draw[flow,omThm] ([yshift=4pt]fcm.west) -- node[above,font=\footnotesize]{its own calls} ([yshift=4pt]f.east);
\draw[flow,omExo,dashed] (eb) -- node[right,font=\footnotesize]{give-up} (fcm);
\end{tikzpicture}

\omcaption{Access to a futures exchange. Membership changes the fee on the
upper path; it does not change the lower one: unless the firm becomes a
\omterm{def:m1:getting-access-futures-options:fcm}{clearing member}, an FCM stands between it and the clearing house and sets its
real margin.}\label{fig:m1:getting-access-futures-options:chain}
\end{omfigure}

\section{Membership: buy, lease or stay outside}

\begin{definition}[Exchange membership and member rate]\label{def:m1:getting-access-futures-options:membership}
An \emph{exchange membership}\index{exchange membership} of a futures exchange
is a transferable right, held by an individual or a firm, that gives access to
the \emph{member rate}\index{member rate}: a schedule of exchange fees lower
than the one applied to everyone else. Memberships are divided by product
group, can be bought and sold at prices the exchange publishes, and can be
leased by the month.
\end{definition}

\begin{dated}{2026-09}{One contract, three prices}\label{dat:m1:getting-access-futures-options:rates}
A futures broker's published comparison gives, for the E-mini S\&P 500, an
exchange fee per side of \$1.18 for a non-member, \$0.47 for the lessee of a
membership and \$0.35 for an owner; for the \omterm{def:m1:index-and-single-stock-futures-worldwide:point}{micro contract}, \$0.20, \$0.07 and
\$0.04. Non-members also pay a regulatory fee of \$0.02 a side. Leasing is a
two-step process: approval as an individual member of the exchange (an
application fee of \$2\,000), then a lease agreement with the holder of a seat.
The exchange has announced changes to its transaction fees effective 1
October 2026; its own fee schedule is the reference.
\end{dated}

\begin{omfigure}
\begin{tikzpicture}
\begin{axis}[width=9.5cm,height=4.8cm,ybar stacked,bar width=18pt,
  ylabel={\$ per side},
  symbolic x coords={non-member,lessee,owner},xtick=data,
  tick label style={font=\small},label style={font=\small},xticklabel style={font=\small\strut},
  ymin=0,ymax=1.7,ymajorgrids,grid style={gray!25},enlarge x limits=0.25,
  legend columns=3,legend style={font=\footnotesize,at={(0.5,-0.2)},anchor=north,draw=none,fill=none,/tikz/every even column/.append style={column sep=8pt}},
  table/col sep=comma]
\addplot[fill=omDef!60,draw=omDef] table[x=status,y=exchange]{figdata/markets-1/30-getting-access-futures-options/per_side.csv};
\addplot[fill=omThm!60,draw=omThm] table[x=status,y=regulatory]{figdata/markets-1/30-getting-access-futures-options/per_side.csv};
\addplot[fill=omExo!50,draw=omExo] table[x=status,y=commission]{figdata/markets-1/30-getting-access-futures-options/per_side.csv};
\legend{exchange fee,regulatory fee,broker commission}
\end{axis}
\end{tikzpicture}

\omcaption{What one side of an E-mini costs: \$1.45, \$0.57 and \$0.45. Exchange
and regulatory fees from \cref{dat:m1:getting-access-futures-options:rates};
the commissions, 25 cents for the outsider and 10 for members, are
illustrative. Data: the tutorial.}\label{fig:m1:getting-access-futures-options:perside}
\end{omfigure}

\begin{proposition}[Break-even volume]\label{prop:m1:getting-access-futures-options:breakeven}
If a route to the exchange costs $c_1$ a side with a fixed monthly cost $K_1$,
and another $c_2 < c_1$ with $K_2 > K_1$, the second is cheaper above
\[
n^* = \frac{K_2 - K_1}{c_1 - c_2}\ \text{sides a month.}
\]
\end{proposition}

With a lease at \$1\,500 a month (an assumption; leases are quoted monthly and
vary by division and with demand) the lessee's 88-cent saving pays for the
lease from 1\,705 sides a month, about 40 round trips a day. Ownership ties up
the price of the membership; at a cost of capital of \$4\,000 a month it beats
leasing from 20\,833 sides a month, a thousand a day
(\cref{fig:m1:getting-access-futures-options:monthly}).

\begin{omfigure}
\begin{tikzpicture}
\begin{axis}[width=11cm,height=5cm,
  xlabel={sides a month (thousand)},ylabel={\$ thousand a month},
  tick label style={font=\small},label style={font=\small},
  xmin=0,xmax=30,ymin=0,ymax=30,grid=major,grid style={gray!25},
  legend columns=3,legend style={font=\footnotesize,at={(0.5,-0.3)},anchor=north,draw=none,fill=none,/tikz/every even column/.append style={column sep=8pt}},
  table/col sep=comma]
\addplot[omThm,very thick] table[x=sides_k,y=non_member_k]{figdata/markets-1/30-getting-access-futures-options/monthly.csv};
\addplot[omDef,very thick,dashed] table[x=sides_k,y=lessee_k]{figdata/markets-1/30-getting-access-futures-options/monthly.csv};
\addplot[black,very thick,dashdotted] table[x=sides_k,y=owner_k]{figdata/markets-1/30-getting-access-futures-options/monthly.csv};
\legend{non-member,lessee,owner}
\end{axis}
\end{tikzpicture}

\omcaption{Monthly cost of trading the E-mini under the three routes. The
lease pays for itself almost at once; ownership only for a firm trading
about a thousand sides a day. Data: the tutorial.}\label{fig:m1:getting-access-futures-options:monthly}
\end{omfigure}

Fees matter in proportion to the tick. A round trip at non-member exchange
rates is 19\% of an E-mini tick and 35\% of a micro tick
(\cref{fig:m1:getting-access-futures-options:ticks}): a strategy that aims to
capture one tick keeps four fifths of it in the large contract and under two
thirds in the small one, before any commission. The smaller contract is the
more expensive, relative to what it can earn, for everyone; and cheap only to
members.

\begin{omfigure}
\begin{tikzpicture}
\begin{axis}[width=10cm,height=4.6cm,ybar,bar width=12pt,
  ylabel={\% of one tick},
  symbolic x coords={non-member,lessee,owner},xtick=data,
  tick label style={font=\small},label style={font=\small},xticklabel style={font=\small\strut},
  ymin=0,ymax=42,ymajorgrids,grid style={gray!25},enlarge x limits=0.25,
  nodes near coords,nodes near coords style={font=\footnotesize,fill=white,inner sep=1pt,/pgf/number format/.cd,fixed,precision=0},
  legend columns=2,legend style={font=\footnotesize,at={(0.5,-0.2)},anchor=north,draw=none,fill=none,/tikz/every even column/.append style={column sep=8pt}},
  table/col sep=comma]
\addplot[fill=omDef!60,draw=omDef] table[x=status,y=es_pct]{figdata/markets-1/30-getting-access-futures-options/ticks.csv};
\addplot[fill=omThm!60,draw=omThm] table[x=status,y=micro_pct]{figdata/markets-1/30-getting-access-futures-options/ticks.csv};
\legend{E-mini (tick \$12.50),micro (tick \$1.25)}
\end{axis}
\end{tikzpicture}

\omcaption{Exchange and regulatory fees of a round trip as a percentage of
one tick. Data: \cref{dat:m1:getting-access-futures-options:rates}.}\label{fig:m1:getting-access-futures-options:ticks}
\end{omfigure}

\section{Incentive and market-maker programmes}

\begin{definition}[Incentive programme]\label{def:m1:getting-access-futures-options:incentive}
An \emph{incentive programme}\index{incentive programme} is a schedule of fee
discounts or rebates that a futures exchange offers, for a product or a
group of products and for a limited period, to participants who meet volume
thresholds or \omterm{def:m1:options-market-structure:mm}{quoting obligations}; membership is often not required.
\end{definition}

Futures exchanges do not pay maker rebates across the board as equity
exchanges do; they target them. New products, products that compete with
another exchange's, back months, and overnight hours receive market-maker
programmes with obligations (maximum width, minimum size, percentage of
time) and stipends or fee waivers in return; high-volume proprietary firms
receive tiered discounts. The terms are filed with the regulator, but
participation is by application, places can be limited, and the programme
can end. A firm whose strategy is profitable only inside a programme has a
business with an expiry date.

\section{Options market-maker appointments}

\begin{definition}[Market-maker appointment]\label{def:m1:getting-access-futures-options:appointment}
A \emph{market-maker appointment}\index{market-maker appointment} is the
registration of a member, on an options exchange, as \omterm{def:m1:what-a-trading-firm-does:market-maker}{market maker} in a set of
option classes. It carries the \omterm{def:m1:options-market-structure:mm}{quoting obligations} and the privileges of
\cref{def:m1:options-market-structure:mm}, and is paid for through trading
permits and appointment fees that grow with the number and the activity of
the classes.
\end{definition}

On an options exchange the fee gap between \omterm{def:m1:what-a-trading-firm-does:market-maker}{market makers} and everyone else
is wider than in futures, and it comes with the allocation entitlements of
\cref{ch:m1:options-market-structure}. The cost side is also larger: permits
on each of a dozen or more exchanges, since the same series trades on all of
them; bulk-quoting bandwidth and ports; the consolidated feed of
\cref{dat:m1:options-market-structure:penny}; and capital for an inventory
of options that cannot be flattened at the close. Options market making has a
minimum efficient scale that futures proprietary trading does not, which is
why it is done by a small number of firms.

\section{Data licences}

The layers are those of \cref{ch:m1:getting-access-equities}: access, display,
non-display, redistribution. Futures add one distinction that catches new
firms: on some exchanges real-time data fees are waived or reduced for
members trading their own account and charged in full to non-members and to
any use that looks like a service to others. The status that lowers the transaction fee
lowers the data bill as well, and the break-even of
\cref{prop:m1:getting-access-futures-options:breakeven} should include both.

\section{Tutorial: when does a lease pay?}\label{tut:m1:getting-access-futures-options:lease}

\begin{tutorial}
\textbf{Goal.} Compute the all-in cost of a side under three routes, the
break-even volumes between them, the effect of a volume incentive, and the
cost of a round trip in ticks. \textbf{End state:} the three data figures of
this chapter.
\begin{tutsteps}
\item \textbf{One side.} Exchange fee and regulatory fee by status, plus the
broker's commission.
\omcode{code/firm/futfees/firm_futfees.py}{10}{29}{Fees by status, routes, and the cost of a month.}\label{lst:m1:getting-access-futures-options:fees}
\item \textbf{Break-even.}
\omcode{code/firm/futfees/firm_futfees.py}{32}{38}{Monthly sides above which the route with the higher fixed cost wins.}\label{lst:m1:getting-access-futures-options:breakeven}
\item \textbf{Incentives} that apply to the whole month once a threshold is
passed: the cliffs of \cref{ch:m1:getting-access-equities} again.
\end{tutsteps}
\textbf{What to change next.} Add the data bill, with a member waiver, to the
fixed cost of each route. Then add a second product in another division and
find whether one membership or two is cheaper for a given split of volume.
\end{tutorial}

\section{Build: the futures fee engine}\label{bld:m1:getting-access-futures-options:futfees}

\begin{build}
\bfield{Purpose} The backtester charges every simulated fill what the firm
would really pay; the firm's finance function decides each quarter whether to
lease, buy or drop memberships.
\bfield{Interface} \texttt{ProductFees(product, exchange\_fee, regulatory\_fee)}
by status; \texttt{Access(status, broker\_commission, fixed\_monthly)};
\texttt{per\_side}, \texttt{monthly\_cost}, \texttt{breakeven\_sides(product,
cheap, dear)}; \texttt{IncentiveTier(min\_sides, discount)} and
\texttt{with\_incentive}; \texttt{cost\_in\_ticks(product, access, tick\_value)}.
\bfield{Rules} Fees are data with effective dates, keyed by product and
status. Fixed costs include the lease or the cost of capital of a purchased
membership. A route that is dearer per side never has a break-even: the
function returns nothing, not a negative number. \omterm{def:m1:futures-contracts-and-exchanges:tick}{Tick values} come from the
contract master (\cref{bld:m1:futures-contracts-and-exchanges:master}).
\bfield{Acceptance tests} \texttt{code/firm/futfees/tests/}: cost per side
and per month; break-even in both directions; an incentive that applies to
the whole month at its threshold; fees against the tick.
\bfield{Stretch} Options on futures and listed options, with fees by
capacity (customer, professional, firm, \omterm{def:m1:what-a-trading-firm-does:market-maker}{market maker}) and by whether the
order added or removed liquidity.
\end{build}

\begin{omsources}
\item National Futures Association, \emph{Futures commission merchant (FCM)
registration}.
\item Optimus Futures, ``CME membership seat lease'' (published comparison of
non-member, lessee and owner exchange fees).
\item CME Group, \emph{Exchange fees for clearing and trading} and
\emph{Membership pricing} (reference pages; not machine-readable at the time
of writing).
\end{omsources}

\section{Exercises}

\begin{exercise}[$\star$]\label{exo:m1:getting-access-futures-options:1}
With the fees of \cref{fig:m1:getting-access-futures-options:perside}, give
the all-in cost of 10\,000 sides in a month for a non-member and for a lessee
paying \$1\,500 for the lease.
\end{exercise}

\begin{exercise}[$\star$]\label{exo:m1:getting-access-futures-options:2}
Compute the two break-even volumes quoted in the text: leasing against
staying outside, and owning (at \$4\,000 a month of capital cost) against
leasing.
\end{exercise}

\begin{exercise}[$\star$]\label{exo:m1:getting-access-futures-options:3}
Express the exchange and regulatory fees of a round trip as a percentage of
one tick for the E-mini (tick \$12.50) and the \omterm{def:m1:index-and-single-stock-futures-worldwide:point}{micro contract} (tick \$1.25), for
each status.
\end{exercise}

\begin{exercise}[$\star\star$]\label{exo:m1:getting-access-futures-options:4}
A strategy makes 400 round trips a day in the E-mini and earns, before fees,
an average of 0.3 tick per round trip. Give its monthly P\&L (21 days) as a
non-member and as a lessee, with the commissions of the figure.
\end{exercise}

\begin{exercise}[$\star\star$]\label{exo:m1:getting-access-futures-options:5}
An exchange offers 5 cents off per side from 50\,000 sides a month and 10
cents from 200\,000, on all sides of the month. A lessee has traded 46\,000
sides with two days left and normally trades 1\,500 a day. What are the last
1\,000 sides worth?
\end{exercise}

\begin{exercise}[$\star\star$]\label{exo:m1:getting-access-futures-options:6}
An FCM asks for 150\% of exchange margin and gives two hours to meet intraday
calls. Another asks for 110\% and ten minutes. For a strategy with large
intraday positions and small overnight ones, which is dearer? What else
would you ask both?
\end{exercise}

\begin{exercise}[$\star\star\star$]\label{exo:m1:getting-access-futures-options:7}
\emph{Coding.} Add a monthly data cost of \$1\,200 for non-members, waived for
lessees and owners, to the routes of the tutorial. Recompute the break-even
between staying outside and leasing. What happens to it if the lease is
\$1\,100?
\end{exercise}

\begin{exercise}[$\star\star\star$]\label{exo:m1:getting-access-futures-options:8}
\emph{Find the flaw.} ``Our market-making strategy in a newly listed future
earns \$180\,000 a month: \$30\,000 from the spread and \$150\,000 from the
exchange's market-maker stipend. We are raising capital to scale it
tenfold.'' What should an investor ask?
\end{exercise}

\section{Problem: Lease or Pay?}

\begin{problem}[{Weekend problem --- a small firm's first membership}]\label{pb:m1:getting-access-futures-options:1}
A two-person firm trades the E-mini through an FCM as a non-member: 600
sides a day, 21 days a month, at the fees of
\cref{fig:m1:getting-access-futures-options:perside}. Its gross trading
profit before fees averages 0.25 tick per side (tick \$12.50). A lease is
offered at \$1\,500 a month; as a lessee its commission falls to 10 cents.

\medskip\noindent\textbf{Part I --- Today.}
\begin{enumerate}
\item Give sides per month and gross profit per month.
\item Give fees per month as a non-member.
\item Give net profit, and fees as a percentage of gross.
\item Give the round trip's cost in ticks.
\end{enumerate}

\medskip\noindent\textbf{Part II --- With a lease.}
\begin{enumerate}[resume]
\item Give the break-even volume in sides per month and per day.
\item Give fees per month as a lessee, lease included.
\item Give the new net profit and the improvement.
\item The application costs \$2\,000 once. How many days of savings repay it?
\item What are the non-financial costs of becoming a member?
\end{enumerate}

\medskip\noindent\textbf{Part III --- Growth.}
\begin{enumerate}[resume]
\item At what volume would owning, at \$4\,000 a month of capital cost, beat
leasing?
\item The firm adds the \omterm{def:m1:index-and-single-stock-futures-worldwide:point}{micro contract}, 2\,000 sides a day at 0.25 tick gross
(tick \$1.25). Give its monthly net as a non-member and as a lessee, with
commissions of 10 and 5 cents.
\item What does question 11 say about which products a non-member can trade
at short horizons?
\item An \omterm{def:m1:getting-access-futures-options:incentive}{incentive programme} offers 10 cents off per side above 20\,000 sides
a month. Give its value to the firm at its E-mini volume.
\item The programme ends next quarter. How should the firm treat its value in
its plans?
\end{enumerate}

\medskip\noindent\textbf{Part IV --- Judgement.}
\begin{enumerate}[resume]
\item Why do exchanges charge non-members three times as much?
\item Who owns memberships and leases them out, and why is that a business?
\item The firm's FCM raises its margin multiple from 120\% to 200\% after a
volatile month. Which matters more to the firm, that or the lease?
\item What would make the firm consider becoming a \omterm{def:m1:getting-access-futures-options:fcm}{clearing member} itself?
\item State the \emph{named result}: the break-even volume of the lease in
sides per month.
\item In one sentence: what does a membership buy?
\end{enumerate}
\end{problem}

\section{Interview questions}

\begin{interviewq}[$\star$ \iqroles{trader, developer}]\label{iq:m1:getting-access-futures-options:1}
What is an FCM and what does it do for a trading firm?
\end{interviewq}

\begin{interviewq}[$\star$ \iqroles{trader, researcher}]\label{iq:m1:getting-access-futures-options:2}
Why do futures exchange fees differ between members and non-members, and by
how much?
\end{interviewq}

\begin{interviewq}[$\star\star$ \iqroles{researcher, trader}]\label{iq:m1:getting-access-futures-options:3}
Your backtest of a one-tick scalping strategy in a \omterm{def:m1:index-and-single-stock-futures-worldwide:point}{micro contract} is
profitable before fees. What do you check next?
\end{interviewq}

\begin{interviewq}[$\star\star$ \iqroles{trader}]\label{iq:m1:getting-access-futures-options:4}
When does it make sense to lease a membership, and when to buy one?
\end{interviewq}

\begin{interviewq}[$\star\star$ \iqroles{trader, bank}]\label{iq:m1:getting-access-futures-options:5}
What would you negotiate with a \omterm{def:m1:getting-access-equities:brokers}{clearing broker}, and what is not
negotiable?
\end{interviewq}

\begin{interviewq}[$\star\star\star$ \iqroles{trader, researcher}]\label{iq:m1:getting-access-futures-options:6}
Why is listed-options market making concentrated in a few firms when
futures proprietary trading is not?
\end{interviewq}
